% !TeX root = ../papers/nc-p3.tex
\section{両切断が欠色を持つ全四元群窓の被覆と正則性判定}
\label{sec:klein-missing-color-coverage}
\begin{theorem}\label{thm:klein-missing-color-coverage}
四位置で四元群の各文字が一回ずつ現れ、最上テンソルは斉次、隣接切断は階数6とする。
両方の切断に少なくとも一つずつ欠色がある場合、印や台への追加仮定なしに次の手順が
標準14641型quadratic Koszul左右AS正則性を必要十分に決定する。
\begin{enumerate}
\item 各切断の色階数が $(0,2,2,2)$ の置換でなければ排除する。
文字原点を選び、総次数および両欠色を零にする。
\item 三つの色ブロックのどれかに零行または零列があれば、二切断のどちらでも指数増大して排除される。
\item それ以外では反復対 $w_{ijij}$ が少なくとも一つ欠ければ常に正則である。
全反復対が非零なら星点で規格化し、四元群印付き交換子の全台判定を用いる。
\end{enumerate}
\end{theorem}
\begin{proof}[Koszul分解の文字別Euler漸化式]
まず一位置のラベルを平行移動して最上テンソルの総次数を零にできる。
反対の二次関係は同じ文字階数を持つ。非自明な指標 $\chi:K\to\{\pm1\}$ について、始点 $p$ の
長さ $n$ の経路の指標値（各文字次元に $\chi$ を掛けた和）を $h_{p,n}$、
関係六空間の値を $s_p$ とする。標準ASなら四項の線型分解はこの多重次数を保ち、一次矢と
三次重なりの各指標値は零、最上の値は1である。$h_{p,0}=1$, $h_{p,1}=0$ とし、
\[
           h_{p,n}=-s_p h_{p+2,n-2}-h_{p+4,n-4}
\]
が成り立つ（負の長さは零）。$s_{p+2}=s_p$ により偶長の母関数は
$1/(1+s_p z+z^2)$ である。各係数の絶対値は標準Hilbert関数以下だから
$|s_p|\le2$ が必要である。実際 $|s_p|>2$ なら相異なる実根の一つが単位円内にあり分子で消えない。
この議論は片方にしか欠色がないときにも各切断に独立に適用できる。

両切断の色を適当に平行移動するには、総次数を零にした後、四位置を
$0,B,B+E,E$ ずらせばよい。欠色を $0$ に取ると、核 $\{0,b\}$ の非自明指標で
$s_p=2r_{p,b}-6$ である。各 $b\ne0$ で $r_{p,b}\ge2$ が必要で、和は6だから全て2となる。
欠色がない側で許される分割は整数の和だけから $(3,1,1,1),(2,2,1,1)$ の二つである。
\end{proof}

\begin{proof}[零行・零列が一方通行だけになることの排除]
自由辺 $i\to j$ を、ある位置で $ij$ がどの関係の台にもない相異なる座標対とする。
欠色0により全ての対角遷移 $ii$ が全位置で使える。従って自由辺に有向閉路があれば
定理\ref{thm:squarefree-window-coverage}の単項式商が指数増大する。
閉路がなければ全自由辺が小から大に向かう共通全順序を選べる。
四頂点の任意の置換は $K$ のアフィン変換なので、共通の付け替えで通常の順序 $0<1<2<3$ にしてよい。

この順序で許される台を二次元直線だけで列挙する。切断色 $b$ の候補は
$w_{i,i+b,i+b+e,i+e}$, $e\ne0$ の十二個の真理値である。行 $i$ と列 $k$ は
それぞれ位置0,2での対 $i(i+b),k(k+b)$ に当たるため、大から小に向かうものは零行・列になれない。
零行と零列を取り除き、残った行・列に二次元空間の直線を一つずつ対応させる。
各行には線型形式の核、各列にはベクトルの像を取る。零成分で結んだ行・列は同じ直線、
非零成分で結んだものは相異なる直線である。零による同値閉包が非零の対を含まず、
行にも列にも二つ以上の同値類があれば、相異なる一般の直線を各類に割り当てて
ちょうどその台の階数2行列を作れる。これで個別の切断での零行・列を含む必要十分条件になる。

十二ビットの4096通りからこの検査で各色171通りを得る。三色を独立に選び、次の切断で
行ラベル・行色・列方向を $(i,b,e)\mapsto(i+b,e,b)$ に置き換える。
この切断でも同じ171通りに入るものは108通りで、どの台も六ブロックの全てに零行も零列も持たない。
\path{computations/verify_klein_missing_color_coverage.py} はこの $171^3$ 通りを二切断のビット位置から直接突き合わせる。
二切断の行列要素の実現を同時に求める必要はなく、必要な台の検査だけで既に零行・列が排除された。
定理\ref{thm:unmarked-klein-pbw}によって最後の段の全次数判定が従う。
\end{proof}
